\section*{Abstract}
There are many Bangla speakers who communicate using regional variations such as Sylhet and Chittagonian. While existing Bangla question answering systems and language models are developed and evaluated on Standard Bangla, this creates a particular gap in the Healthcare sector and education sector in the domain where reliable information is important. This study proposes a dialect aware multi domain Bangla QA framework based on router driven mixture of expert architecture. These frameworks collect questions in Bangla dialect or standard Bangla and responses according to that input. A light weight routing model first recognises standard Bangla and linguistic variety of the input and its domain. These inputs and questions are then normalized into standard Bangla if it is recognized by the model as a dialect before being routed to the proper domain so that each domain is supported by a dedicated retrieval augmented generation (RAG) pipeline so that it matches the response generation in relevant external evidence. There will be multiple predetermined language models where the framework compares several open source Bangla adapted TigerLLM and other suitable architectures. These models are adapted using efficient fine tuning methods such as LoRA and QLoRA. The framework is determined using Bangla question answering and domain specific resources including NCTB-QA, BanglaRQA and BanglaCHQ. The evaluation covers dialect identification and normalization, retrieval effectiveness, end-to-end QA performance, clarification and abstention behavior. Then the results are separately reported across that specific dialect language or standard Bangla with component wise studies and comparisons across candidate base models. This study aims to reduce the performance disparity between standard Bangla and dialect users by developing a modular retrieval, grounded and domain hour QA framework that provides standard Bangla responses while being computationally practical for low resource Bangla NLP.

\vspace{0.5cm}
\textbf{Keywords:} Bangla NLP, Bengali medical chatbot, Bengali educational chatbot, Chittagonian, Sylheti, dialect normalization, dialect routing, question answering, retrieval-augmented generation, medical safety, Gemma 4, TigerLLM, LoRA, QLoRA, mixture of adapters, NCTB-QA, BanglaRQA, BanglaCHQ, low-resource NLP
\clearpage
